\chapter{Haseo Ki and the Thorin programme}\label{chap:ki}\chaptermark{Ki and the Thorin programme}

Haseo Ki and his coauthors have studied the objects this book computes with: Fourier transforms of rapidly decaying positive kernels, their zeros under the heat flow and under repeated differentiation, and zeta functions that share the functional equation of $\zeta$ but not its Euler product. The de Bruijn Newman flow of Section~\ref{ch:dbn} and the Thorin condition of Chapter~\ref{chap:foundations} meet in this work, and this chapter reads it through the Thorin dictionary: a zero of an even entire function is an atom of the Thorin measure of its reciprocal, and the question whether the reciprocal is a generalised gamma convolution (GGC) is the question whether the zeros are real.

Four results come out. The theorem of Ki, Kim and Lee on the de Bruijn Newman flow says that after any positive time the reciprocal of the flowed function is a GGC transform up to a finite rational factor, and de Bruijn's strip theorem confines the atoms of that factor to confocal parabolas that shrink with time (Section~\ref{ki:sec:flow}). Read together with the theorem of Rodgers and Tao, the flow gives an obstruction: no property that survives a Gaussian tilt of the P\'olya density can imply RH (Proposition~\ref{ki:prop:obstruction}). Ki's theorem on repeated differentiation is a statement about moment tilts of the P\'olya density, and it implies that under RH the reciprocals of the derivatives are GGC transforms whose clocks converge in law to the exit time of Brownian motion from an interval (Section~\ref{ki:sec:deriv}); a conjecture, supported by computation, gives the rate at which the atoms converge. Ki's approximation theorems for Epstein type zeta functions and the Davenport Heilbronn functions indicate that a Thorin argument for RH must use the Euler product (Section~\ref{ki:sec:approx}). Ki's work on the zeros of $\zeta'$ concerns the spacing of the zeros (Section~\ref{ch:spacing}) and gives no route to positivity; his uniqueness theorems in the extended Selberg class and his work in descriptive set theory are not used.

Status words follow the rest of the book. A statement is proved when a proof or a pointer is given, computed when a script of the book produced it, and heuristic, conditional or open as stated. Results of Ki and his coauthors are cited, not reproved.
